\lecture{27}{Матрица Лапласа и спектральная связность}

\sourcevideo
    {https://www.youtube.com/playlist?list=PLOzRYVm0a65fhKDS8cYIC7YCuVGJ4FOJx}
    {Week 8: Lecture 27: Basics of Graph Theory-2}

Пусть неориентированная взвешенная коммуникационная сеть задана
графом
\[
    \mathcal G=(\mathcal V,\mathcal E),
    \qquad
    \mathcal V=\{1,\ldots,N\}.
\]
Её матрица смежности равна
\[
    A=(a_{ij})_{i,j=1}^{N},
    \qquad
    a_{ij}=a_{ji}\ge0,
\]
а матрица степеней имеет вид
\[
    D=\operatorname{diag}(d_1,\ldots,d_N),
    \qquad
    d_i=\sum_{j=1}^{N}a_{ij}.
\]
Матрица Лапласа объединяет локальную структуру рёбер и глобальные
свойства связности. Её спектр определяет скорость консенсусных и
распределённых алгоритмов.

\section{Лапласиан и его квадратичная форма}

\begin{definition}[Матрица Лапласа]
Матрицей Лапласа неориентированного взвешенного графа называется
\[
    L=D-A.
\]
\end{definition}

Её элементы равны
\[
    L_{ij}
    =
    \begin{cases}
        d_i, & i=j,\\
        -a_{ij}, & i\ne j.
    \end{cases}
\]
Для простого невзвешенного графа ненулевые внедиагональные элементы
равны \(-1\). Поскольку \(A=A^{\mathsf T}\), матрица \(L\)
симметрична. Сумма элементов каждой строки равна нулю:
\[
    \sum_{j=1}^{N}L_{ij}
    =
    d_i-\sum_{j=1}^{N}a_{ij}
    =
    0.
\]
Следовательно,
\[
    L\mathbf1=0,
    \qquad
    \mathbf1^{\mathsf T}L=0,
\]
где \(\mathbf1\in\mathbb R^N\) --- вектор из единиц. Поэтому нуль
всегда является собственным значением Лапласиана.

Для произвольного \(x\in\mathbb R^N\)
\[
\begin{aligned}
    (Lx)_i
    &=
    d_ix_i-\sum_{j=1}^{N}a_{ij}x_j
    \\
    &=
    \sum_{j=1}^{N}a_{ij}(x_i-x_j).
\end{aligned}
\]
Таким образом, компоненту \((Lx)_i\) можно вычислить по состоянию
вершины \(i\) и состояниям её соседей.

Квадратичная форма имеет вид
\[
\begin{aligned}
    x^{\mathsf T}Lx
    &=
    \frac12
    \sum_{i,j=1}^{N}
    a_{ij}(x_i-x_j)^2
    \\
    &=
    \sum_{\{i,j\}\in\mathcal E}
    a_{ij}(x_i-x_j)^2,
\end{aligned}
\]
где во второй сумме каждое неориентированное ребро учитывается один
раз. Все слагаемые неотрицательны.

\begin{theorem}
Матрица Лапласа неориентированного графа положительно
полуопределена:
\[
    L\succeq0.
\]
\end{theorem}

Следовательно, все собственные значения \(L\) вещественны и
неотрицательны.

\section{Ядро и компоненты связности}

Из формулы для квадратичной формы следует
\[
    x^{\mathsf T}Lx=0
\]
тогда и только тогда, когда
\[
    x_i=x_j
\]
на каждом ребре положительного веса. Поскольку \(L\succeq0\), это
эквивалентно условию \(Lx=0\). Значит, любой вектор из ядра
Лапласиана постоянен на каждой компоненте связности.

Если граф связен, то
\[
    \ker L=\operatorname{span}\{\mathbf1\}.
\]
Если граф состоит из \(r\) компонент
\[
    \mathcal V_1,\ldots,\mathcal V_r,
\]
то
\[
    \ker L
    =
    \operatorname{span}
    \left\{
        \mathbf1_{\mathcal V_1},
        \ldots,
        \mathbf1_{\mathcal V_r}
    \right\},
\]
где \(\mathbf1_{\mathcal V_s}\) --- индикатор компоненты
\(\mathcal V_s\). Поэтому
\[
    \dim\ker L=r.
\]

\begin{theorem}
Кратность нулевого собственного значения Лапласиана равна числу
компонент связности графа.
\end{theorem}

После подходящей перенумерации вершин матрицы смежности, степеней и
Лапласа принимают блочно-диагональный вид:
\[
    L=
    \begin{pmatrix}
        L_1&0&\cdots&0\\
        0&L_2&\cdots&0\\
        \vdots&\vdots&\ddots&\vdots\\
        0&0&\cdots&L_r
    \end{pmatrix}.
\]
Каждый блок соответствует одной связной компоненте и имеет ровно одно
нулевое собственное значение.

\section{Алгебраическая связность и консенсус}

Упорядочим собственные значения:
\[
    0=\lambda_1(L)
    \le
    \lambda_2(L)
    \le
    \cdots
    \le
    \lambda_N(L).
\]
Граф связен тогда и только тогда, когда нулевое собственное значение
простое:
\[
    \mathcal G\text{ связен}
    \quad\Longleftrightarrow\quad
    \lambda_2(L)>0.
\]

\begin{definition}[Алгебраическая связность]
Второе наименьшее собственное значение
\[
    \lambda_2(L)
\]
называется алгебраической связностью, или собственным значением
Фидлера.
\end{definition}

Для симметричного Лапласиана справедлива вариационная формула Рэлея:
\[
    \lambda_2(L)
    =
    \min_{\substack{x\ne0\\x^{\mathsf T}\mathbf1=0}}
    \frac{x^{\mathsf T}Lx}{x^{\mathsf T}x}.
\]
Следовательно, для каждого \(x\perp\mathbf1\)
\[
    x^{\mathsf T}Lx
    \ge
    \lambda_2(L)\lVert x\rVert^2.
\]

Для произвольного \(x\in\mathbb R^N\) определим среднее
\[
    \overline x
    =
    \frac1N\mathbf1^{\mathsf T}x
\]
и рассогласование
\[
    x_\perp=x-\overline x\,\mathbf1.
\]
Тогда \(x_\perp\perp\mathbf1\), а из \(L\mathbf1=0\) следует
\[
    Lx=Lx_\perp,
    \qquad
    x^{\mathsf T}Lx
    =
    x_\perp^{\mathsf T}Lx_\perp.
\]
Для связного графа получаем основную оценку
\[
    x^{\mathsf T}Lx
    \ge
    \lambda_2(L)
    \lVert
        x-\overline x\,\mathbf1
    \rVert^2.
\]
Она показывает, насколько сумма локальных различий на рёбрах
контролирует глобальное расстояние до консенсуса. Малое
\(\lambda_2(L)\) указывает на слабую связь между частями сети, а
большое значение означает более сильное спектральное перемешивание.

Если каждый агент хранит вектор \(x_i\in\mathbb R^d\), то для
составного состояния
\[
    \mathbf x=
    \begin{pmatrix}
        x_1\\
        \vdots\\
        x_N
    \end{pmatrix}
    \in\mathbb R^{Nd}
\]
используется оператор \(L\otimes I_d\). При связном графе
\[
    \ker(L\otimes I_d)
    =
    \{
        \mathbf1\otimes v:
        v\in\mathbb R^d
    \},
\]
поэтому
\[
    (L\otimes I_d)\mathbf x=0
    \quad\Longleftrightarrow\quad
    x_1=\cdots=x_N.
\]

\section{Нормализованные Лапласианы}

Матрицу
\[
    L=D-A
\]
называют ненормализованным Лапласианом. Если все степени положительны,
определяются симметричный нормализованный Лапласиан
\[
    L_{\mathrm{sym}}
    =
    D^{-1/2}LD^{-1/2}
    =
    I-D^{-1/2}AD^{-1/2}
\]
и Лапласиан случайного блуждания
\[
    L_{\mathrm{rw}}
    =
    D^{-1}L
    =
    I-D^{-1}A.
\]
Они связаны преобразованием подобия
\[
    L_{\mathrm{rw}}
    =
    D^{-1/2}
    L_{\mathrm{sym}}
    D^{1/2}
\]
и потому имеют одинаковые собственные значения.

Ненормализованный Лапласиан учитывает абсолютные веса рёбер, поэтому
вершины большой степени сильнее влияют на квадратичную форму.
Нормализация сравнивает связи относительно степеней вершин. По этой
причине \(L_{\mathrm{sym}}\) и \(L_{\mathrm{rw}}\) используются в
спектральной кластеризации и моделях случайного блуждания.

При наличии изолированных вершин обычные матрицы \(D^{-1}\) и
\(D^{-1/2}\) не определены. Тогда нормализация задаётся только на
вершинах положительной степени или используется соглашение с
псевдообратной матрицей.

\section{Расстояние и диаметр графа}

\begin{definition}[Графовое расстояние]
Расстоянием между вершинами \(i\) и \(j\) связного графа называется
длина кратчайшего пути между ними:
\[
    \operatorname{dist}_{\mathcal G}(i,j).
\]
\end{definition}

\begin{definition}[Диаметр]
Диаметром связного графа называется величина
\[
    \operatorname{diam}(\mathcal G)
    =
    \max_{i,j\in\mathcal V}
    \operatorname{dist}_{\mathcal G}(i,j).
\]
\end{definition}

Для пути \(P_N\)
\[
    \operatorname{diam}(P_N)=N-1,
\]
а для звезды с \(N\ge3\)
\[
    \operatorname{diam}(\mathcal G)=2.
\]
Диаметр определяет минимальное число последовательных передач,
необходимых для доставки одного сообщения между наиболее удалёнными
вершинами.

Малый диаметр не означает, что стандартный итерационный алгоритм
усреднения достигает точного консенсуса за такое же число итераций.
Скорость численного алгоритма зависит также от весов и спектра
матрицы обновления. Диаметр является комбинаторной характеристикой, а
\(\lambda_2(L)\) --- спектральной. У хорошо связанных графов диаметр
обычно мал, а алгебраическая связность относительно велика, однако
универсального равенства между этими величинами нет.

\section{Динамика консенсуса и локальная реализация}

Рассмотрим непрерывную систему
\[
    \dot x=-Lx.
\]
Среднее состояние сохраняется, поскольку
\[
\begin{aligned}
    \frac{\mathrm d}{\mathrm dt}
    \left(
        \frac1N\mathbf1^{\mathsf T}x
    \right)
    &=
    -\frac1N\mathbf1^{\mathsf T}Lx
    =
    0.
\end{aligned}
\]
Следовательно, рассогласование
\[
    x_\perp=x-\overline x\,\mathbf1
\]
удовлетворяет уравнению
\[
    \dot x_\perp=-Lx_\perp.
\]
Для функции Ляпунова
\[
    V(x_\perp)
    =
    \frac12\lVert x_\perp\rVert^2
\]
получаем
\[
\begin{aligned}
    \dot V
    &=
    -x_\perp^{\mathsf T}Lx_\perp
    \\
    &\le
    -\lambda_2(L)
    \lVert x_\perp\rVert^2
    \\
    &=
    -2\lambda_2(L)V.
\end{aligned}
\]
Поэтому для связного графа
\[
    \lVert x_\perp(t)\rVert
    \le
    e^{-\lambda_2(L)t}
    \lVert x_\perp(0)\rVert.
\]
Собственное значение Фидлера тем самым непосредственно задаёт
гарантированную экспоненциальную скорость непрерывного консенсуса.

Явно хранить плотную матрицу \(L\) не требуется. Каждый агент
вычисляет
\[
    (Lx)_i
    =
    \sum_{j\in\mathcal N_i}
    a_{ij}(x_i-x_j)
\]
по сообщениям соседей. Для разреженного графа стоимость вычисления
всех компонент имеет порядок
\[
    O(\lvert\mathcal E\rvert),
\]
а для состояний размерности \(d\) ---
\[
    O(\lvert\mathcal E\rvert d).
\]
Именно сочетание локальной реализуемости и спектрального контроля
глобального рассогласования делает Лапласиан основным оператором
децентрализованной оптимизации.
